problem:  
Let \((V^{n}, g, J)\) be a Hermitian vector space and let \(\mathcal{R}\) denote the real vector space of **algebraic Kähler curvature tensors**, that is, \(U(n)\)-invariant tensors satisfying the usual Kähler symmetries.  The curvature operator \(R:\Lambda^{2}V\to \Lambda^{2}V\) preserves the real subspace of \((1,1)\)-forms, denoted \(\Lambda^{1,1}(V)\).  Equip \(\Lambda^{1,1}(V)\) with the inner product induced by \(g\).  Let \(\Lambda^{1,1}_0(V)\subset \Lambda^{1,1}(V)\) denote the primitive (traceless) forms, on which \(U(n)\) acts irreducibly.  Define  
\[
R^{\#}:\Lambda^{1,1}_0(V)\longrightarrow \Lambda^{1,1}_0(V), \qquad
\langle R^{\#}(\alpha),\beta\rangle = \langle R(\alpha),\beta\rangle,
\]
as the curvature operator restricted to \(\Lambda^{1,1}_0(V)\).

Call a form \(\alpha\in\Lambda^{1,1}_0(V)\) **simple** if \(\alpha=X\wedge JY-JX\wedge Y\) for some orthonormal pair \((X,Y)\).  Simple primitive \((1,1)\)-forms are the Kähler analogues of isotropic 2‑vectors in the Riemannian setting.

Define the **complex isotropic–positive curvature cone**  
\[
\mathcal{C}_{\text{CIP}}
   =\{\,R\in\mathcal{R}\mid 
       \langle R^{\#}(\alpha),\alpha\rangle\ge0 
       \text{ for all simple }\alpha\in\Lambda^{1,1}_0(V)\,\}.
\]
It is immediate that \(\mathcal{C}_{\text{OBC}}\subseteq \mathcal{C}_{\text{CIP}}\subseteq \mathcal{C}_{\text{Ric}}\), but whether these inclusions are strict or convexity holds globally is unknown.

**Open problem:**  
(1)  Determine whether, for every compact Kähler manifold \((M^{2n},g,J)\) satisfying \(\operatorname{Rm}(M)\in\mathcal{C}_{\text{CIP}}\) pointwise, the **Frankel intersection property** holds:
\[
X^p, Y^q\subset M\ \text{closed complex submanifolds with}\ p+q\ge n
\quad\Rightarrow\quad
X\cap Y\neq\varnothing.
\]
(2)  If not universally true, characterize the **maximal \(U(n)\)-invariant Frankel-effective subcone**
\(\mathcal{C}_{\text{CIP}}^{\text{Frankel}}\subseteq\mathcal{C}_{\text{CIP}}\)
for which the intersection property still holds.  
(3)  Investigate whether the cone \(\mathcal{C}_{\text{CIP}}\) or its Frankel-effective part is preserved under the Kähler–Ricci flow.

Why is it a "good" problem:  
This formulation now uses a rigorously defined and intrinsically \(U(n)\)-invariant cone built from the curvature operator on \(\Lambda^{1,1}_0(V)\), eliminating ad hoc coordinate expressions.  It sits naturally within the known representation‑theoretic framework of Kähler curvature cones and links several active research directions:

* **Curvature–Topology Frontier:** It aims to identify the minimal algebraic positivity ensuring global intersection rigidity, extending Frankel’s phenomenon beyond holomorphic bisectional curvature.
* **Representation Theory and Convex Geometry:** The cone \(\mathcal{C}_{\text{CIP}}\) is defined in a representation‑theoretic manner, allowing explicit algebraic exploration of positivity domains in curvature space.
* **Flow Invariance and Geometric Analysis:** Asking whether this cone is preserved under Kähler–Ricci flow connects complex differential geometry, analytic evolution equations, and convex‑cone dynamics.
* **Conceptual Economy:** The problem’s statement uses only an elementary curvature positivity notion and the classical intersection property, yet resolving it would require deep interplay between curvature algebra, geodesic geometry, and global topology.

Because it refines the notion of intermediate curvature positivity in a mathematically coherent, structurally founded way and connects curvature algebra to global intersection phenomena and Ricci‑flow dynamics, this problem is concise, profound, and genuinely situated at a research frontier—a “good” mathematical problem.